% Tinjauan pustaka
\section{Tinjauan Pustaka}

\subsection{Struktur Tes UTBK-SNBT 2026}
Setiap subtes UTBK-SNBT 2026 dilaporkan pada skala standar dengan rerata populasi $M = 500$. Khusus LBI, sertifikat 2026 memuat dua nilai sebagai tanggapan atas anggapan bahwa peserta berlatar saintek lebih diuntungkan ketika memilih program studi soshum \autocite{detik2026lbi}. Statistik deskriptif univariat nasional yang dipakai dalam paper ini berasal dari parameter nasional tidak resmi (belum terverifikasi). Asal dan keterbatasannya diuraikan pada Bagian Metode.

\subsection{Korelasi Antarkemampuan Kognitif (\textit{Positive Manifold})}
\textcite{spearman1904} pertama kali mendokumentasikan \textit{positive manifold}, fenomena mendasar dalam kajian kemampuan kognitif. Analisis ulang atas lebih dari 460 kumpulan data \autocite{carroll1993} melahirkan \textit{Carroll's three-stratum theory}, yang menempatkan \textit{g-factor} pada strata tertinggi. Tinjauan komprehensif \textcite{jensen1998} semakin memperkuat landasan teoretis eksistensi \textit{g-factor}. Besaran rerata korelasi antarsubtes yang tersirat dari model semacam ini bergantung pada komposisi rangkaian tes (\textit{test battery}) dan karakteristik sampel. Oleh karena itu, penulis tidak menetapkan rentang acuan numerik apriori.

Selain rerata korelasi tunggal, kerangka faktor kelompok dalam \textit{Carroll's three-stratum theory} \autocite{carroll1993} menyiratkan bahwa subtes kognitif berpotensi membentuk dimensi laten (kemampuan yang tidak diukur langsung) yang lebih spesifik di bawah \textit{g-factor}. Contohnya pembedaan antara kemampuan verbal dan kuantitatif. Pengujian formal terhadap struktur multifaktor semacam ini lazim dilakukan melalui CFA \autocite{brown2015cfa}. CFA membandingkan kecocokan statistik antara \textit{g-factor} tunggal dan model berjenjang yang memuat faktor kelompok spesifik.

\subsection{Penelitian Terdahulu dan Celah Penelitian}
Berbagai bukti empiris berskala besar mendukung korelasi positif antarsubtes tes kognitif. \textcite{jintelligence2026} menghimpun puluhan ribu koefisien korelasi dari ratusan rangkaian tes kognitif dan akademik terstandardisasi. Hasil itu sejalan dengan meta-analisis struktural \textcite{consensusmodel2026} pada domain hubungan kognitif-membaca. Pada tingkat instrumen tunggal, \textcite{waisiv2013} melaporkan korelasi antarsubtes yang moderat pada \textit{Wechsler Adult Intelligence Scale} (WAIS-IV).

Sebagai kasus pembanding yang relevan, studi validitas nasional College Board \autocite{westrick2019} melibatkan 223.858 mahasiswa di 171 perguruan tinggi empat tahun. Studi itu melaporkan korelasi antara skor seksi \textit{Evidence-Based Reading and Writing} dan Matematika pada SAT, yang strukturnya menyerupai komponen literasi dan kuantitatif pada UTBK-SNBT. Korelasi tersebut sebesar $0,58$ pada \textit{raw sample} dan $0,78$ setelah dikoreksi untuk \textit{restriction of range} multivariat terhadap populasi pengikut SAT nasional. Selisih kedua nilai itu muncul pada sampel mahasiswa yang sudah terseleksi dan menunjukkan bahwa \textit{restriction of range} menekan korelasi yang teramati secara substansial.

Sejauh penelusuran pustaka yang dilakukan, belum ditemukan publikasi ilmiah maupun laporan resmi yang menyajikan struktur korelasi antarsubtes UTBK-SNBT secara spesifik. Dengan demikian, paper ini diposisikan untuk mengisi celah informasi empiris itu.

\subsection{Kerangka Berpikir}
Paper ini bersifat eksploratif-deskriptif, sehingga hipotesis statistik formal tidak dirumuskan. Istilah konfirmatori pada CFA merujuk pada teknik pembandingan model faktor. Sintesis literatur psikometri di atas mengarah pada satu ekspektasi kualitatif. Apabila UTBK-SNBT memiliki karakteristik pengukuran yang serupa dengan rangkaian tes kognitif terstandardisasi lainnya, korelasi antarsubtesnya diperkirakan positif secara konsisten. Apabila korelasi pada \textit{raw sample} ternyata relatif lemah dan disertai indikasi \textit{restriction of range} pada parameter sampel teramati (sebagaimana dipaparkan pada Bagian Hasil), kondisi itu diposisikan sebagai indikasi awal bias metodologis dalam pengambilan sampel. Kondisi itu tidak ditafsirkan sebagai cerminan struktur psikometri UTBK-SNBT yang berbeda secara mendasar dari tes kognitif pada umumnya.

Kerangka faktor kelompok pada \textit{Carroll's three-stratum theory} juga mengarah pada ekspektasi struktur faktor yang berjenjang. Dalam struktur itu, domain Literasi (verbal) dan Numerasi (kuantitatif) menjadi dua faktor kelompok di bawah \textit{g-factor}. CFA membandingkan ekspektasi tersebut dengan \textit{g-factor} tunggal.